\documentclass[10pt,a4paper]{article}
\usepackage{amsmath,amssymb,amsthm,graphicx}
\usepackage{amsmath,amssymb,amsthm,hyperref,xfrac}
\providecommand{\ie}{i.e.\ }
\providecommand{\skpt}{\smallskip}
\usepackage[scr=rsfs,cal=euler]{mathalfa}
\multlinegap0pt

\hyphenation{begin decom-po-si-tion sequence}
\newenvironment{enumeratea}
{\bgroup\def\theenumi{\alph{enumi}}\begin{enumerate}}
{\end{enumerate}\egroup}

\newcommand{\ba}{{\boldsymbol a}}
\newcommand{\bb}{{\boldsymbol b}}
\newcommand{\bd}{{\boldsymbol d}}
\newcommand{\bu}{{\boldsymbol u}}
\newcommand{\bx}{{\boldsymbol x}}
\newcommand{\bmy}{{\boldsymbol y}}

\newcommand{\bD}{{\boldsymbol D}}
\newcommand{\bL}{{\boldsymbol L}}
\newcommand{\bP}{{\boldsymbol P}}
\newcommand{\bV}{{\boldsymbol V}}
\newcommand{\bX}{{\boldsymbol X}}
\newcommand{\bZ}{{\boldsymbol Z}}

\newcommand{\Psfrac}[2]{(\sfrac{#1}{#2})}
\newcommand{\psfrac}[2]{\sfrac{(#1)}{#2}}
\newcommand{\bigO}[1]{O(#1)}

\usepackage{mathtools}
\usepackage{bbm}

\newcommand{\NN}{\mathbb{N}}
\newcommand{\ZZ}{\mathbb{Z}}
\newcommand{\RR}{\mathbb{R}}
\newcommand{\CC}{\mathbb{C}}
\newcommand{\PP}{\mathbb{P}}
\newcommand{\EE}{\mathbb{E}}

\newcommand{\boldalpha}{\boldsymbol{\alpha}}
\newcommand{\bolddelta}{\boldsymbol{\delta}}
\newcommand{\boldrho}{\boldsymbol{\rho}}
\newcommand{\dee}{\,\mathrm{d}}
\newcommand{\one}{\mathbbm{1}}
\newcommand{\BC}{\mathcal{B}}
\newcommand{\EC}{\mathcal{E}}
\newcommand{\RS}{\mathscr{R}}
\newcommand{\prob}[1]{\PP[ #1 ]}
\newcommand{\bigoo}[1]{O\!\left(#1 \right)}
\newcommand{\abs}[1]{\left\lvert #1 \right\rvert}
\newcommand{\floor}[1]{\left\lfloor #1 \right\rfloor}
\newcommand{\riemint}[4]{\int_{#1}^{#2}\! #3 \, \mathrm{d}#4}
\newcommand{\maxnorm}[1]{\left\| #1 \right\|_{\infty}}

\renewcommand{\Re}{\textup{Re}}
\renewcommand{\Im}{\textup{Im}}

\newcommand{\dtv}{\mathrm{d}_{\mathrm{TV}}}
\newcommand{\dw}{\mathrm{d}_{\mathrm W}}

\let\epsilon\varepsilon

\newcommand{\pextra}{P_{\mathtt{extra}}}
\newcommand{\parratia}{P_{\mathtt{Arratia}}}

\newcommand{\simplex}{\Delta^{k-1}}

\newcommand{\bg}{\big}
\newcommand{\bgg}{\Big}
\newcommand{\bggg}{\bigg}

\newcommand{\Rnt}{R_{\mathrm{NT}}}
\newcommand{\Rprob}{R_{\mathrm{PD}}}
\newcommand{\ECnt}{\EC_{\mathrm{NT}}}
\newcommand{\ECprob}{\EC_{\mathrm{PD}}}

\newcommand{\Bnt}{\mathcal{B}_{\mathrm{NT}}}
\newcommand{\Bprob}{\mathcal{B}_{\mathrm{PD}}}

\DeclareMathOperator{\vard}{Var}
\DeclareMathOperator{\Dir}{Dir}

\theoremstyle{plain}
\newtheorem{cor}{Corollary}[section]
\newtheorem{lem}[cor]{Lemma}
\newtheorem{prop}[cor]{Proposition}
\newtheorem{conj}[cor]{Conjecture}
\newtheorem{theo}{Theorem}

\theoremstyle{definition}
\newtheorem{example}{Example}
\newtheorem{deff}[cor]{Definition}
\newtheorem{rem}{Remark}[section]
\newtheorem*{rems}{Remark}
\newtheorem*{remss}{Remarks}

\makeatletter
\ams@newcommand{\multiint}[1]{\DOTSI\protect\MultiIntegral{#1}}
\renewcommand{\MultiIntegral}[1]{%
\edef\ints@c{\noexpand\intop
\ifnum#1=\z@\noexpand\intdots@\else\noexpand\intkern@\fi
\replicate{#1-2}{\noexpand\intop\noexpand\intkern@}%
\noexpand\intop
\noexpand\ilimits@
}%
\futurelet\@let@token\ints@a
}
\makeatother
\ExplSyntaxOn
\cs_new:Npn \replicate #1 #2 { \prg_replicate:nn { #1 } { #2 } }
\ExplSyntaxOff
\title{The law of the empirical frequency vector differs from the Dirichlet limit only off two exceptional events}
\author{Tony Haddad, Dimitris Koukoulopoulos}
\date{doi:10.5802/jep.317; CC BY 4.0}
\begin{document}
\maketitle
\noindent\textit{Editorial note.} The counted statement is the author's labelled lem environment at source0.tex lines 1266-1272; its proof environment opens at line 1274, the line immediately after the blank separator, so the pairing is by adjacency and the retained proof range 1274-1298 is the whole environment.\par\medskip
\noindent\textit{Editorial note.} Two environments are retained as context because the counted proof cites them: the author's coupling lemma (source lines 1159-1166), used at lines 1274 and 1295, and the numbered transition estimate eq:transition delta to delta* (source lines 1219-1222), used at line 1287.\par\medskip
\noindent\textit{Editorial note.} This article carries no label-number dump, so no counter values are imposed: the retained environments print with the numbering the wrapper produces, and the identity of the counted item is its author label lem:sym diff D and Z.\par\medskip
\noindent\textit{Editorial note.} The retained preamble is the publisher layout adapter, the native font and hyphenation packages and the native mathematical macro and theorem block of the pinned source; the article bibliography is bundled because the retained proof carries no citation of its own but the source does elsewhere in the same file.\par\medskip
\noindent\textit{Editorial note.} r002 note: the two numbered equations eq:def of Theta\_x (source lines 504-507) and eq:defn delta (source lines 1153-1156) are retained as context because the retained coupling-lemma context refers to them at source line 1165; r001 failed only on those two undefined references.\par\medskip
\noindent\textit{Editorial scope.} One original result of the pinned publisher article, counted once: the comparison of the law of the empirical frequency vector with its Dirichlet limit off two exceptional events (author label lem:sym diff D and Z).\par\medskip
\noindent\textit{Allowed standard prerequisites (editorial).}\begin{itemize}
\item Empirical frequency vector and Dirichlet limit. delta\_\{f,x\} is the vector of normalised logarithms of the ordered parts of a random factorisation, and Z is the Dirichlet limit vector.. They are the two random vectors whose laws the statement compares.
\item The events E, E\_cnt(A) and E\_prob(A). E is the coupling event of the author's coupling lemma; E\_cnt(A) and E\_prob(A) are the events that the two relevant closed balls avoid the boundary of the measurable set A.. The statement bounds the discrepancy by the complements of these events plus 1/log x, and the proof works on their intersection.
\item Coupling distance bound. On the coupling event E one has a uniform bound on the maximal norm of the difference of delta\_\{f,x\} and Z, in terms of log(x/N\_x), log s(N\_x) and Theta\_x.. It supplies the radius within which a point of the connecting segment is exhibited.
\end{itemize}\medskip
\begin{equation}
\label{eq:def of Theta_x}
\Theta_x \coloneqq \sum_{i \ge 1} r(V_i\log x).
\end{equation}

\begin{equation}
\label{eq:defn delta}
\bolddelta_{f, x} \coloneqq \bgg(\frac{\log D_{f,x,1}}{\log N_x},\dots,\frac{\log D_{f,x,k}}{\log N_x}\bgg)
\end{equation}

\begin{lem}
\label{lem:coupling for factorizations}
Let $x\ge2$ and assume the above notation. There exists a coupling of $N_x$, $\bV$, $\bD_{f, x}$ and $\bZ$ such that $\PP[\mathcal E^c] \ll \sfrac{1}{\log x}$ with
\[
\mathcal E \coloneqq \bgg\{\maxnorm{\bolddelta_{f, x} - \bZ} \le \frac{2 \cdot \log(x/N_x) + 3 \cdot \log s(N_x) + 2\cdot \Theta_x}{\log x}\bgg\},
\]
with $\bolddelta_{f, x}$ defined as in \eqref{eq:defn delta}, $\Theta_x$ defined as in \eqref{eq:def of Theta_x}, and $\maxnorm \cdot$ being the supremum norm in $\RR^k$.
\end{lem}

\begin{equation}
\label{eq:transition delta to delta*}
\maxnorm{ \bolddelta_{f, x} - \bolddelta_{f, x}^*} \le \frac{\log(x/N_x)}{\log x}.
\end{equation}

\begin{lem}
\label{lem:sym diff D and Z}
For any measurable set $A\subseteq\simplex$, we~have
\[
\PP\bg[\bolddelta_{f, x} \in A\bg] = \PP\bg[\bZ \in A\bg] + O\Bigl(\PP\bg[\ECnt(A)^c\bg] + \PP\bg[\ECprob(A)^c\bg] + \frac{1}{\log x}\Bigr).
\]
\end{lem}

\begin{proof} Let $\EC$ be as in Lemma \ref{lem:coupling for factorizations}. Hence, it suffices to show that
\begin{equation}
\label{eq:sym diff D and Z on E}
\bgg| \PP\bg[\EC\cap\{\bolddelta_{f, x} \in A\}\bg] - \PP\bg[\EC\cap\{\bZ \in A\}\bg]\bgg| \le \PP\bg[\ECnt(A)^c\bg] + \PP\bg[\ECprob(A)^c\bg].
\end{equation}
We shall prove the following stronger statement: if $\EC\cap \ECnt(A)\cap \ECprob(A)$ occurs, then $\bolddelta_{f, x}\in A$ if, and only if, $\bZ\in A$.

Indeed, assume that $\EC\cap \ECnt(A)\cap \ECprob(A)$ occurs, and we also have $\bolddelta_{f, x}\in A$. Aiming for a contradiction, we~also assume that $\bZ\notin A$.

Let
\[
\Rnt' \coloneqq (\log x)^{-1}(2\cdot \log(x/N_x) + 3\cdot \log s(N_x)).
\]
Note that $B(\bolddelta_{f, x}, \Rnt') \subseteq B(\bolddelta_{f, x}^*, \Rnt)$ by \eqref{eq:transition delta to delta*}. By assumption of $\ECnt(A)$, we~know that \hbox{$B(\bolddelta_{f, x},\Rnt')\cap\partial A=\emptyset$}. Since $\bolddelta_{f, x} \in A$ and $B(\bolddelta_{f, x},\Rnt')\cap \simplex$ is connected (being the intersection of two convex sets), we~have that
\[
B(\bolddelta_{f, x},\Rnt')\cap \simplex \subseteq A.
\]
And by assumption of $\ECprob(A)$, we~have $\bZ \in \simplex \setminus A$ and $B(\bZ,\Rprob)\cap\partial A=\emptyset$, which implies
\[
B(\bZ,\Rprob)\cap \simplex \subseteq \simplex \setminus A.
\]
In particular, we~find that $B(\bolddelta_{f, x}, \Rnt')\cap B(\bZ,\Rprob)\cap \simplex =\emptyset$. On the other hand, we~know that $\maxnorm{\bolddelta_{f, x}-\bZ}\le \Rnt' +\Rprob$ when $\EC$ occurs by Lemma \ref{lem:coupling for factorizations}. In~particular, there exists a point $\bmy$ on the line segment connecting $\bolddelta_{f, x}$ and $\bZ$ such that $\maxnorm{\bolddelta_{f, x}-\bmy}\le \Rnt'$ and $\maxnorm{\bmy-\bZ}\le \Rprob$. Since the sets $\simplex$, $B(\bolddelta_{f, x},\Rnt')$ and $B(\bZ,\Rprob)$ are convex, we~conclude that $\bmy$ lies in their intersection. But we had seen before that this intersection is the empty set. We~have thus arrived at a contradiction. This completes the proof that if $\EC\cap \ECnt(A)\cap \ECprob(A)$ occurs, and we also know that $\bolddelta_{f, x}\in A$, then we must also have that $\bZ\in A$.

Conversely, we~may show by a simple variation of the above argument that if $\EC\cap \ECnt(A)\cap \ECprob(A)$ occurs, and we also know that $\bZ\in A$, then we must also have that $\bolddelta_{f, x}\in A$. This completes the proof of the lemma.
\end{proof}

\begin{thebibliography}{99}
\bibitem{Arratia98} R. Arratia - “On the central role of scale invariant Poisson processes on (0,∞)”, in Microsurveys in discrete probability (Princeton, NJ, 1997), DIMACS Ser. Discrete Math. Theoret. Comput. Sci., vol. 41, American Mathematical Society, Providence, RI, 1998, p. 21-41
\bibitem{Arratia02} R. Arratia - “On the amount of dependence in the prime factorization of a uniform random integer”, in Contemporary combinatorics, Bolyai Soc. Math. Stud., vol. 10, János Bolyai Math. Soc., Budapest, 2002, p. 29-91
\bibitem{ABT06} R. Arratia, A. D. Barbour \& S. Tavaré - “A tale of three couplings: Poisson-Dirichlet and GEM approximations for random permutations”, Combin. Probab. Comput. 15 (2006) no. 1-2, p. 31-62
\bibitem{BareikisManstavicius07} G. Bareikis \& E. Manstavičius - “On the DDT theorem”, Acta Arith. 126 (2007) no. 2, p. 155-168
\bibitem{Bill72} P. Billingsley - “On the distribution of large prime divisors”, Period. Math. Hungar. 2 (1972), p. 283-289
\bibitem{delaBretecheTenenbaum16} R. de la Bretèche \& G. Tenenbaum - “Sur les processus arithmétiques liés aux diviseurs”, Adv. in Appl. Probab. 48 (2016) no. A, p. 63-76
\bibitem{DDT79} J.-M. Deshouillers, F. Dress \& G. Tenenbaum - “Lois de répartition des diviseurs. I”, Acta Arith. 34 (1979) no. 4, p. 273-285 (loose errata)
\bibitem{DonnellyTavare87} P. Donnelly \& S. Tavaré - “The population genealogy of the infinitely-many neutral alleles model”, J. Math. Biol. 25 (1987) no. 4, p. 381-391
\bibitem{Feng10} S. Feng - The Poisson-Dirichlet distribution and related topics. Models and asymptotic behaviors, Probability and its Appl., Springer, Heidelberg, 2010
\bibitem{JKB} N. L. Johnson, S. Kotz \& N. Balakrishnan - Discrete multivariate distributions, Wiley Series in Probab. and Stat.: Applied Probab. and Stat., John Wiley \& Sons, Inc., New York, 1997
\bibitem{Kingman93} J. F. C. Kingman - Poisson processes, Oxford Studies in Probab., vol. 3, The Clarendon Press, Oxford University Press, New York, 1993
\bibitem{KotzBal00} S. Kotz, N. Balakrishnan \& N. L. Johnson - Continuous multivariate distributions. Vol. 1. Models and applications, Wiley Series in Probab. and Stat.: Applied Probab. and Stat., Wiley-Interscience, New York, 2000
\bibitem{Koukoulo19} D. Koukoulopoulos - The distribution of prime numbers, Graduate Studies in Math., vol. 203, American Mathematical Society, Providence, RI, 2019
\bibitem{LastPenrose18} G. Last \& M. Penrose - Lectures on the Poisson process, Institute of Math. Stat. Textbooks, vol. 7, Cambridge University Press, Cambridge, 2018
\bibitem{Ken22} S.-K. Leung - “Dirichlet law for factorisation of integers, polynomials and permutations”, Math. Proc. Cambridge Philos. Soc. 175 (2023) no. 3, p. 649-676
\bibitem{KenWeb24} S.-K. Leung - “Dirichlet law with arbitrary parameters for factorizations of integers” (2025), Retrieved on 24 August 2025 from https://sites.google.com/view/sunkaileung/miscellaneous
\bibitem{MV07} H. L. Montgomery \& R. C. Vaughan - Multiplicative number theory. I. Classical theory, Cambridge Studies in Advanced Math., vol. 97, Cambridge University Press, Cambridge, 2007
\bibitem{NyandwiSmati13} S. Nyandwi \& A. Smati - “Distribution laws of pairs of divisors”, Integers 13 (2013), article ID A13, 13 pages
\bibitem{Tenenbaum00} G. Tenenbaum - “A rate estimate in Billingsley’s theorem for the size distribution of large prime factors”, Q. J. Math. 51 (2000) no. 3, p. 385-403
\end{thebibliography}
\end{document}
